\documentclass[a4paper,12pt]{article}

\usepackage[utf8]{inputenc}
\usepackage[T1]{fontenc}
\usepackage{amsmath}
\usepackage{booktabs}
\usepackage{graphicx}
\usepackage{geometry}
\geometry{margin=2.5cm}
\usepackage{xcolor}
\usepackage{listings}
\usepackage{hyperref}
\usepackage{adjustbox}
\hypersetup{colorlinks=true, linkcolor=blue!50!black, urlcolor=blue!50!black}
\renewcommand{\tablename}{Tabela}
\renewcommand{\lstlistingname}{Código}

% ---------- Configuração do realce de sintaxe (C) ----------
\definecolor{codebg}{RGB}{247,247,247}
\definecolor{codekw}{RGB}{0,0,180}
\definecolor{codecomment}{RGB}{0,128,0}
\definecolor{codestring}{RGB}{163,21,21}
\definecolor{framecolor}{RGB}{200,200,200}

\lstdefinestyle{cstyle}{
    language=C,
    backgroundcolor=\color{codebg},
    basicstyle=\ttfamily\footnotesize,
    keywordstyle=\color{codekw}\bfseries,
    commentstyle=\color{codecomment}\itshape,
    stringstyle=\color{codestring},
    numberstyle=\tiny\color{gray},
    numbers=left,
    numbersep=6pt,
    frame=single,
    rulecolor=\color{framecolor},
    breaklines=true,
    breakatwhitespace=false,
    showstringspaces=false,
    tabsize=4,
    captionpos=b,
    columns=flexible,
    morekeywords={size_t},
    extendedchars=true,
    literate=
        {á}{{\'a}}1 {à}{{\`a}}1 {â}{{\^a}}1 {ã}{{\~a}}1
        {é}{{\'e}}1 {ê}{{\^e}}1 {í}{{\'i}}1
        {ó}{{\'o}}1 {ô}{{\^o}}1 {õ}{{\~o}}1
        {ú}{{\'u}}1 {ü}{{\"u}}1 {ç}{{\c c}}1
        {Á}{{\'A}}1 {À}{{\`A}}1 {Â}{{\^A}}1 {Ã}{{\~A}}1
        {É}{{\'E}}1 {Ê}{{\^E}}1 {Í}{{\'I}}1
        {Ó}{{\'O}}1 {Ô}{{\^O}}1 {Õ}{{\~O}}1
        {Ú}{{\'U}}1 {Ü}{{\"U}}1 {Ç}{{\c C}}1
}

\title{Tarefa 8: Coerência de Cache e Falso Compartilhamento}
\author{Israel Soares de Castro Filho \\ Matrícula: 20250070780}
\date{}

\begin{document}

\maketitle

\section{Introdução}

Este relatório analisa quatro versões de um programa em C/OpenMP que estima $\pi$ pelo método de Monte
Carlo: sorteiam-se pontos $(x,y)$ no quadrado $[0,1]\times[0,1]$ e verifica-se quantos caem dentro do
quarto de círculo unitário ($x^2+y^2 \le 1$); a razão entre acertos e total de pontos converge para
$\pi/4$.

O objetivo é comparar duas formas de acumular os acertos de cada thread -- variável privada somada a um
total global via \texttt{\#pragma omp critical}, ou vetor compartilhado com uma posição por thread e soma
serial ao final -- combinadas com dois geradores de números aleatórios: \texttt{rand()} (estado global,
não thread-safe por padrão) e um gerador com estado privado por thread. A análise explica os tempos
observados com base em coerência de cache e falso compartilhamento (\textit{false sharing}).

\section{Metodologia}

O programa implementa quatro variantes, todas com o laço de amostragem distribuído via
\texttt{\#pragma omp for}:

\begin{enumerate}
    \item \textbf{V1 -- \texttt{rand()} + critical}: contador local por thread, somado ao total global em
    um bloco \texttt{critical} ao final da região paralela.
    \item \textbf{V2 -- \texttt{rand()} + vetor}: cada thread grava seu contador local em \texttt{acertos\_por
    \_thread[id]}; 
    a soma é feita em laço serial após a região paralela.
    \item \textbf{V3 -- Xorshift32 + critical}: igual a V1, trocando \texttt{rand()} por um gerador
    Xorshift32 com semente privada por thread.
    \item \textbf{V4 -- Xorshift32 + vetor}: igual a V2, também com Xorshift32.
\end{enumerate}

Como o ambiente é Windows e \texttt{rand\_r()} é uma função POSIX indisponível no MinGW
/MSVC, ela foi
substituída por um Xorshift32 com estado (semente) puramente privado por thread -- preservando a
característica central que se quer testar: um gerador sem estado global compartilhado.

O programa foi executado com $8$ threads, para $n\_pontos$ de $10^5$ a $10^9$ (passos $\times 10$),
medindo o tempo de cada versão com \texttt{omp\_get\_wtime()}. O código completo está no
Apêndice~\ref{sec:codigo}.

\section{Resultados}

\begin{lstlisting}[style=cstyle, language=, numbers=none, caption={Saída da execução com 8 threads}, label={lst:saida}]
===============================================================
        ESTIMATIVA DE PI - METODO DE MONTE CARLO
===============================================================
PI_REAL = 3.141592653589793
Threads = 8
Intervalo de n_pontos = 100000 ate 1000000000

N_PONTOS     VERSAO   PI_MEDIDO    ERRO_ABS     ERRO_REL(%)  TEMPO(s)
--------------------------------------------------------------------------
100000       1        3.11740000   0.02419265   0.770076     0.0020
100000       2        3.13556000   0.00603265   0.192025     0.0010
100000       3        3.13220000   0.00939265   0.298977     0.0000
100000       4        3.13220000   0.00939265   0.298977     0.0000
--------------------------------------------------------------------------
1000000      1        3.13700000   0.00459265   0.146189     0.0060
1000000      2        3.14625600   0.00466335   0.148439     0.0050
1000000      3        3.14120800   0.00038465   0.012244     0.0010
1000000      4        3.14120800   0.00038465   0.012244     0.0000
--------------------------------------------------------------------------
10000000     1        3.14031280   0.00127985   0.040739     0.0990
10000000     2        3.13953440   0.00205825   0.065516     0.0710
10000000     3        3.14231040   0.00071775   0.022847     0.0100
10000000     4        3.14231040   0.00071775   0.022847     0.0090
--------------------------------------------------------------------------
100000000    1        3.14113016   0.00046249   0.014722     0.4030
100000000    2        3.14198412   0.00039147   0.012461     0.4460
100000000    3        3.14182792   0.00023527   0.007489     0.0890
100000000    4        3.14182792   0.00023527   0.007489     0.0830
--------------------------------------------------------------------------
1000000000   1        3.14143575   0.00015691   0.004994     4.0620
1000000000   2        3.14160874   0.00001609   0.000512     4.1530
1000000000   3        3.14157712   0.00001554   0.000495     0.7930
1000000000   4        3.14157712   0.00001554   0.000495     0.7610
--------------------------------------------------------------------------
\end{lstlisting}

\begin{table}[h]
    \centering
    \caption{Tempo de execução (s) por versão e por tamanho de amostra}
    \label{tab:tempos}
    \begin{adjustbox}{width=\textwidth}
        \begin{tabular}{@{}lcccc@{}}
            \toprule
            \textbf{N\_PONTOS} & \textbf{V1: rand+critical} & \textbf{V2: rand+vetor} & \textbf{V3: xorshift+critical} & \textbf{V4: xorshift+vetor} \\
            \midrule
            $10^5$  & 0,0020 & 0,0010 & 0,0000 & 0,0000 \\
            $10^6$  & 0,0060 & 0,0050 & 0,0010 & 0,0000 \\
            $10^7$  & 0,0990 & 0,0710 & 0,0100 & 0,0090 \\
            $10^8$  & 0,4030 & 0,4460 & 0,0890 & 0,0830 \\
            $10^9$  & 4,0620 & 4,1530 & 0,7930 & 0,7610 \\
            \bottomrule
        \end{tabular}
    \end{adjustbox}
\end{table}

O erro relativo cai com o aumento de $n\_pontos$ em todas as versões, como esperado de um estimador cuja
incerteza decresce com $1/\sqrt{n}$. As quatro versões convergem para o mesmo $\pi$; a diferença relevante
está nos tempos de execução.

\section{Análise}

\textbf{rand() vs. Xorshift32.} Esse é o fator que mais pesa nos tempos. Para
$n\_pontos = 10^9$, V1 e V2 (\texttt{rand()}) levam cerca de $4,1\,\text{s}$,
contra apenas $0,76$--$0,79\,\text{s}$ de V3 e V4 (Xorshift32) -- uma diferença
de $\sim 5,3\times$. A razão é que \texttt{rand()} mantém um único estado
global, compartilhado por todas as threads, \emph{sem nenhuma proteção
explícita} contra acesso concorrente -- as chamadas simultâneas a
\texttt{rand()} em V1 e V2 constituem, a rigor, uma condição de corrida,
 já que múltiplas threads leem e escrevem a mesma variável
de estado sem sincronização. Como o laço chama \texttt{rand()} duas vezes por
iteração, esse estado é lido/escrito centenas de milhões de vezes, e cada
escrita obriga o protocolo de coerência de cache a invalidar e retransmitir a
mesma linha de cache entre os núcleos (um caso de compartilhamento verdadeiro).

Isso serializa boa parte do trabalho que deveria ser paralelo -- não porque
exista uma seção crítica formal protegendo o estado, mas porque o hardware
precisa manter a coerência da linha de cache a cada acesso concorrente. Vale
notar que essa condição de corrida em V1 e V2 é, tecnicamente, comportamento
indefinido pela linguagem C; os resultados de $\pi$ obtidos continuam
estatisticamente razoáveis porque o método de Monte Carlo tolera ruído nos
números gerados, mas isso não deve ser confundido com um padrão de acesso
seguro. O Xorshift32, por outro lado, guarda sua semente como variável
totalmente privada (registrador/pilha de cada thread), sem estado
compartilhado, e por isso escala quase linearmente com o número de threads.

\textbf{\texttt{critical} vs. vetor compartilhado.} Comparando V1/V2 e V3/V4, as diferenças são pequenas
(ex.: $0,793$s vs.\ $0,761$s em V3/V4 para $10^9$). Isso ocorre porque, em ambas as versões, o contador é
acumulado localmente (\texttt{acertos\_local}, em registrador) durante todo o laço de $n\_pontos$
iterações; a sincronização -- seção crítica ou escrita no vetor -- acontece apenas \emph{uma vez por
thread}, ao final. Como esse custo é $O(1)$ por thread e não $O(n\_pontos)$, ele é irrelevante frente ao
custo de gerar e testar centenas de milhões de pontos, e por isso a escolha entre as duas estratégias
pouco altera o tempo total.

\textbf{Por que quase não há falso compartilhamento.} O vetor \texttt{acertos\_por\_thread} é de
\texttt{long long} (8 bytes); como uma linha de cache tem 64 bytes, até 8 contadores de threads vizinhas
compartilham a mesma linha. Se cada thread escrevesse em \texttt{acertos\_por
\_thread[id]} \emph{a cada
iteração} do laço, haveria falso compartilhamento severo (invalidações constantes da linha entre threads
que não compartilham dado algum, apenas a linha de cache). Mas no código essa escrita ocorre só
\emph{uma vez}, após o laço paralelo -- então o número de possíveis invalidações é da ordem do número de
threads ($8$), não de $n\_pontos$ (até $10^9$). É por isso que V2 e V4 não são mais lentas que V1 e V3:
o padrão de acesso ao vetor é raro demais para que o falso compartilhamento apareça no tempo medido.

\textbf{Resumo.} Para $10^9$ pontos: V4 ($0,761$s) $<$ V3 ($0,793$s) $\ll$ V1 ($4,062$s) $<$ V2
($4,153$s). Os dois níveis de desempenho, bem separados, mostram que o gerador aleatório (estado global vs.\
privado) é o fator dominante; a estratégia de acumulação (\texttt{critical} vs.\ vetor) é secundária.

\section{Conclusão}

O gargalo de desempenho não está na acumulação dos resultados parciais, mas na
geração dos números aleatórios. Um gerador com estado global (\texttt{rand()}),
acessado concorrentemente sem qualquer proteção, sofre uma condição de corrida
que gera tráfego intenso de coerência de cache entre os núcleos e serializa
boa parte do programa. Um gerador com estado privado por thread elimina
esse problema e escala quase linearmente. 

Já \texttt{critical} e vetor
compartilhado se mostraram equivalentes, pois ambos custam apenas uma operação
por thread, e não uma por iteração -- inclusive o vetor, apesar de
compartilhar linhas de cache entre threads, não sofre falso compartilhamento
relevante por ser acessado tão raramente. A lição geral é que o estado
verdadeiramente compartilhado (aqui, o gerador de números aleatórios) é o que
determina a escalabilidade, muito mais do que pequenas seções críticas ou
vetores usados com moderação.

\newpage
\appendix

\section{Código-Fonte}
\label{sec:codigo}

O código-fonte completo utilizado nesta atividade está disponível abaixo.

\lstinputlisting[style=cstyle, caption={Código-fonte completo (\texttt{pi\_monte\_carlo\_omp.c})}]{../src/pi_monte_carlo_omp.c}

\end{document}